\label{sec:other-methods}

\subsection{Webster}

Webster's priority for the $(s+1)$-th seat is $p_i/(s+0.5) = 2p_i/(2s+1)$.
In exact arithmetic:

\begin{verbatim}
fn webster_priority(pop: u128, seats: u32) -> Priority {
    Priority {
        numerator: 4 * pop * pop,
        denominator: (2 * seats as u128 + 1) * (2 * seats as u128 + 1),
    }
}
\end{verbatim}

The factor of 4 in the numerator cancels the factor of 4 implicit
in squaring the denominator's $2s+1$: the comparison
$p_i/(s_i+0.5) > p_j/(s_j+0.5)$ becomes
$p_i(s_j+0.5) > p_j(s_i+0.5)$, equivalently
$2p_i(2s_j+1) > 2p_j(2s_i+1)$, equivalently
$(2p_i)^2(2s_j+1)^2 > (2p_j)^2(2s_i+1)^2$,
which is the cross-multiplication form.

Overflow check: maximum value is $(2 \times 3.95 \times 10^7)^2 \times
(2 \times 52+1)^2 = (7.9 \times 10^7)^2 \times 105^2
\approx 6.9 \times 10^{17}$, within \texttt{u128}.

\subsection{Adams}

Adams uses $P_i^{\text{Ad}}(s) = p_i/s$:

\begin{verbatim}
fn adams_priority(pop: u128, seats: u32) -> Priority {
    Priority {
        numerator: pop * pop,
        denominator: (seats as u128) * (seats as u128),
    }
}
\end{verbatim}

The comparison $p_i/s_i > p_j/s_j$ becomes $p_i \times s_j > p_j \times s_i$
(or equivalently $p_i^2 s_j^2 > p_j^2 s_i^2$).

\subsection{Jefferson}

Jefferson uses $P_i^{\text{Jeff}}(s) = p_i/(s+1)$
(starting from $s=0$ for the first seat):

\begin{verbatim}
fn jefferson_priority(pop: u128, seats: u32) -> Priority {
    Priority {
        numerator: pop * pop,
        denominator: (seats as u128 + 1) * (seats as u128 + 1),
    }
}
\end{verbatim}

Jefferson initialises with 0 seats per state (not 1), allocating
all $H$ seats via priority queue, and then applies the constitutional
floor as post-processing.

\subsection{Hamilton}

Hamilton does not use the priority queue (see Section~\ref{sec:architecture}).
The key implementation detail is representing exact quotas as rational
fractions to avoid floating-point precision loss in the fractional-part comparison:
the fractional part $f_i = q_i - \lfloor q_i \rfloor$ is stored as
$(p_i \times H \mod N, N)$ in integer arithmetic.

\subsection{Edge Cases}

All methods handle:
\begin{itemize}
  \item \textbf{Single state}: Returns $H$ seats (trivially).
  \item \textbf{Exact quotas}: If $q_i$ is an integer for all $i$
    (and $\sum q_i = H$), all methods agree.
  \item \textbf{Constitutional floor}: States with exact quota $< 1$
    receive 1 seat; the divisor adjusts to compensate.
  \item \textbf{Maximum state count}: Tested with $n = 50$, $H = 435$
    (the federal case) and with $n = 100$, $H = 1000$ (stress test).
\end{itemize}
